\section{Controlled Transfer Tests}\label{sec:controlled-placement}\label{sec:tau-cover-bridge}
The experiments operationalize the audit's three bindings: observation, execution and matched evaluation. E5--E7 isolate mechanisms using saved outputs and reference programs; E9 evaluates real LLM-generated trajectories with training-only calibration. Seven author-written judgments were recorded before E5--E7; their reassessment is retained in Supplement S12. These cases test mechanisms, not the audit's external decision accuracy.

For E6/E7/E9, a cell meets the \emph{evaluated-suite criterion} when maximum BA has an upper 95\% bound below 0.60 and maximum AUROC is below 0.65. A lower BA bound above 0.60 indicates above-threshold leakage; other cells are inconclusive. We recompute the suite maximum within each of 2,000 common bootstrap draws: tasks for E6/E7 and customers for E9, retaining their tasks and repeats together. Intervals condition on fitted predictions and do not cover unknown attackers.

A byte-transparent relay captures post-handshake TLS application records on loopback connections. TLS~1.3 session tickets are disabled. E5--E8 use saved text or stand-in planners; E9 invokes a real model and executes its structured actions in the local public retail simulator. No production account or tool is modified.

\subsection{Transport: buffered and streaming emitters}
\paragraph{Buffered emitter}
The corpus has one prompted brief and one prompted expanded \texttt{gpt-6-luna} response for each of 24 topics (12 training and 12 test topics); these are designed conditions, not natural length samples. In a two-link TLS~1.2 replay, a frontend that decodes and re-chunks the upstream text leaves the terminal link 100\% separable when only the upstream link is padded; padding both links reduces the evaluated classifier to 50\%. This reproduces NetEcho's mechanism \cite{netecho2025}. With workload-matched denominators, 512-byte buckets cost 2.30$\times$ and 1.43$\times$ for two sentence-prefix workloads (T2, T3); a cap fitted on training lengths (1,280 bytes) yields 50\% at 2.18$\times$ for verbosity (T1). Under offsets, 1024-byte buckets and random tails, refitting ($D_2$) gains 13.5--50.0 points over $D_0$; constant caps give no gain. Supplement S11 retains all cells.

\paragraph{E5: the same texts as a token stream}
These figures describe an emitter that frames a response only after generation is complete. LLM APIs usually stream, and the streaming attacks in our corpus observe per-token records \cite{weiss2024prompt,whisperleak2025}. E5 emits the same 48 texts as server-sent chunk events, one TLS record per token, over TLS~1.2 (AES-128-GCM) and TLS~1.3 (AES-256-GCM). Token boundaries use o200k\_base as a proxy; the generating model's tokenizer and timing were not saved, so latency is reported in token intervals. The framings are per-event random padding of 1--64 bytes (an obfuscation field), a constant 144-byte event, 8-token batches padded to a constant size, and constant events plus dummy events up to the largest training token count. Besides $D_0$ and $D_2$ threshold and sequence classifiers, a closed-world attacker profiles all 48 candidate responses in an independent run and matches each observation to the nearest record-length sequence.

Table~\ref{tab:transfer-streaming} resolves P1--P3. Per-event padding hides token lengths but not their number. Verbosity remains 100\% separable for both attackers, and identification falls only from 100\% to 58--62\% (chance 2.1\%). Batching sends fewer envelopes (0.21$\times$ bytes) and drives $D_0$ to 50\%, while $D_2$ still scores 100\%. A constant-shape stream reaches 50\% at 2.42$\times$ stream bytes and doubles completion time (208 versus 105 token intervals); two test outputs overflow the fitted count and remain identifiable. The buffered cap's 2.19$\times$ is a different cost object: its first byte arrives only after the full response. The 2.18$\times$ result therefore does not price a streaming deployment; the result illustrates the emitter and cost bindings in Check~1 and R9.

\begin{table}[t]\centering\footnotesize
\caption{E5: the same 48 saved responses emitted by a buffered or a token-streaming frontend (TLS~1.3 loopback captures). Bytes are relative to the unpadded emitter of the same kind. TTFB/Done: first and last delivered byte in token intervals (accounting proxy). Accuracies on 24 held-out outputs (12 topics); ID: closed-world top-1 identification among 48 responses (chance 2.1\%).}\label{tab:transfer-streaming}
\setlength{\tabcolsep}{1.6pt}
\begin{tabularx}{\columnwidth}{@{}Yrrrrrrr@{}}\toprule
 & & & & \multicolumn{2}{c}{T1 acc.} & T3 & \\\cmidrule(lr){5-6}\cmidrule(lr){7-7}
Emitter / framing & Bytes$\times$ & TTFB & Done & $D_0$ & $D_2$ & $D_2$ & ID\\\midrule
Buffered, unpadded & 1.00 & 105 & 105 & 100.0 & 100.0 & 91.7 & 87.5\\
Buffered, training cap & 2.19 & 105 & 105 & 50.0 & 50.0 & 50.0 & 2.1\\
Stream, per token & 1.00 & 1 & 105 & 100.0 & 100.0 & 87.5 & 100.0\\
+ random event pad & 1.29 & 1 & 105 & 100.0 & 100.0 & 87.9 & 62.3\\
+ constant event size & 1.23 & 1 & 105 & 100.0 & 100.0 & 87.5 & 58.3\\
8-token batches, const. & 0.21 & 8 & 105 & 50.0 & 100.0 & 87.5 & 20.8\\
Const.\ events + count cap & 2.42 & 1 & 208 & 50.0 & 50.0 & 50.0 & 6.3\\
\bottomrule\end{tabularx}
\par\smallskip\raggedright
$D_2$ is the best of total-byte threshold, record-count threshold and a logistic sequence model, each fitted on defended training outputs. Random padding: 10 socket-measured draws plus an independent profiling run. TLS~1.2 gives identical accuracies and byte ratios within 0.03. Tokens: o200k\_base proxy. Count cap: 208 (T1) and 88 (T3) events fitted on training outputs; 2/24 and 5/24 test outputs overflow it.
\end{table}


\subsection{Execution: state scope and randomized cover}
\paragraph{State scope}
We replay the 115 public retail reference programs. Same-item writes change workflow state: the item-modification tool sets the status to \texttt{pending (item modified)} and appends a refund entry when item and price are unchanged (173/291 pre-insertion units). A no-change requirement $f(s,a)=s$ thus differs from repeated-call idempotence $f(f(s,a),a)=f(s,a)$. Cancelling another simulated user's order preserves the active account in every unit while changing the shared database in every unit. One cover call at a known first or last position is removed by a position or count rule with 100\% accuracy. These are properties of one simulator and of a deliberately weak cover design (Supplement S11).

\paragraph{E6: randomized cover against policy-aware attackers}
E6 asks what a properly randomized cover achieves in the same tool-dispatch view, which shows ordered tool names and response lengths, including failures. Cover calls are drawn from other users' reference programs, inserted at uniformly random positions, and executed in a per-session sandbox (shadow) or, as a control, live. The policies add one cancellation with probability $0.5$ or $1$, pad the cancellation count to two, interleave one or two complete decoy programs, or equalize the program \emph{shape}. The last pads every tool's count to its maximum over the public program population, a label-free but full-population template. Decoy pools and framing caps also use all 115 tasks. This is known-population defense calibration; attacker holdout does not make it unseen-task defense evaluation. Attackers are gradient boosting and logistic regression over tool counts, bigrams and response lengths, fitted on clean ($D_0$) or defended ($D_2$) sessions, plus a count rule. We use task-grouped 5-fold cross-validation over all 115 tasks and three seeds, with two targets: cancellation (18 tasks) and return (32 tasks).

The dispatch columns of Table~\ref{tab:transfer-cover} show that adding one cancellation always takes $D_0$ to 50\% BA, but $D_2$ recovers 100\% (P2). Equalizing the cancellation count defeats the count rule (50\%), whereas a learned model using correlated tools reaches 76.9\% (P4). Decoy programs, the read-only defense pattern of WRA \cite{jeong2026wra}, leave 83.5--91.9\% because most decoy programs contain no cancellation. Shape equalization gives 50.0\% and 54.5\% point estimates, at 6.7 cover calls per real call. Sandboxed cover keeps database equality at 100\%; the same cancellation executed live yields 0\%. Equality verifies state preservation under the isolation design, not task quality, authorization or deployment utility.

\begin{table*}[t]\centering\footnotesize
\caption{E6/E7: randomized cover on 115 retail reference programs $\times$ 3 seeds. Entries are the best held-out balanced accuracy (\%) over all attackers for the cancellation (C, 18 positive tasks) and return (R, 32) targets. Dispatch view: tool names and response lengths. Network view: TLS~1.2 record lengths on the agent--tool link under raw, 1024-byte bucket or constant framing.}\label{tab:transfer-cover}
\setlength{\tabcolsep}{3.5pt}
\begin{tabularx}{\textwidth}{@{}Yrrrrrrrrrrr@{}}\toprule
 & \multicolumn{4}{c}{Dispatch view (E6)} & \multicolumn{7}{c}{Network view (E7)}\\\cmidrule(lr){2-5}\cmidrule(lr){6-12}
 & & & & & \multicolumn{2}{c}{Raw} & \multicolumn{2}{c}{Bucket} & \multicolumn{2}{c}{Constant} & \\\cmidrule(lr){6-7}\cmidrule(lr){8-9}\cmidrule(lr){10-11}
Cover policy & C & R & Calls+ & DB eq. & C & R & C & R & C & R & Bytes$\times$\\\midrule
No cover & 100.0 & 100.0 & 0.00 & 100.0 & 100.0 & 95.7 & 60.1 & 73.9 & 59.7 & 53.6 & 3.4\\
Cancel, $p{=}0.5$ & 86.1 & 100.0 & 0.10 & 100.0 & --- & --- & --- & --- & --- & --- & ---\\
Cancel, $p{=}1$ & 100.0 & 100.0 & 0.20 & 100.0 & --- & --- & --- & --- & --- & --- & ---\\
Cancel, $p{=}1$, live execution & 100.0 & 100.0 & 0.20 & 0.0 & --- & --- & --- & --- & --- & --- & ---\\
Cancel count $\to 2$ & 76.9 & 100.0 & 0.35 & 100.0 & 72.7 & 95.7 & 67.5 & 66.0 & 64.9 & 59.1 & 4.6\\
One decoy program & 91.9 & 85.9 & 1.01 & 100.0 & 91.9 & 85.9 & 60.0 & 59.1 & 59.7 & 56.8 & 6.9\\
Two decoy programs & 83.5 & 78.1 & 1.98 & 100.0 & --- & --- & --- & --- & --- & --- & ---\\
Shape equalization & 50.0 & 54.5 & 6.71 & 100.0 & 51.2 & 60.5 & 55.3 & 50.6 & 54.4 & 50.0 & 26.3\\
\bottomrule\end{tabularx}
\par\smallskip\raggedright
Calls+: cover calls per real call. DB eq.: orders/users equal to reference. Cover is sampled from other users' reference programs and placed at random positions; all policies except ``live'' run cover in a per-session sandbox. Attackers: $D_0$/$D_2$ gradient boosting and logistic regression, a count rule and, on the network view, a size-profiling decoder; task-grouped 5-fold cross-validation. Bytes$\times$: constant framing relative to raw framing without cover. Suite-maximum intervals: Supplement S12. Defense templates/pools/caps use all 115 tasks; attacker fitting uses training folds. No network configuration meets the criterion for both targets.
\end{table*}


\subsection{E7: composing execution and transport on the agent--tool link}
Figure~\ref{fig:system-model} draws observation and execution as separate paths; a remote tool server joins them on one encrypted connection. E7 replays four cover policies over a TLS~1.2 agent--tool link with one length-prefixed request and response per call (4,140 sessions). Framing is raw, 1024-byte buckets or constant (512-byte requests; 3,072-byte responses, the 95th percentile, with larger responses padded to 1024-byte multiples). The network observer sees direction-tagged record lengths and no tool names. In addition to the learned models, a profiling attacker learns a size-to-tool table by calling the public tool API on training-fold calls, decodes test traces and applies the dispatch attacks.

Under raw framing, request-only identification is 89.9\% with cancellation-stratified folds and 91.4\% with return-stratified folds (582 calls; training-majority baseline 29.4\%). The former 95.7\% figure is in-sample size-table purity, not held-out accuracy. The downstream decoder additionally uses response sizes and is fitted on training-fold calls. Count and decoy cover leave 72.7--95.7\% BA on the network link (P5). They already fail in the dispatch view, so this is persistence of leakage, not a protected-source-to-unprotected-target transition.

Framing alone leaves bucketed return inference at 73.9\%; other framing-only cells are inconclusive. Shape cover with raw framing has cancellation BA 51.2\% but return BA 60.5\%; with constant framing the values are 54.4\% and 50.0\%. Thus several point estimates are near chance. No network configuration meets the evaluated-suite criterion for \emph{both} targets. The joint constant-framing configuration costs 26.3$\times$ bytes and 7.7$\times$ calls, with database equality 100\%; this is one measured operating point, not a minimum necessary cost (P7). Supplement S12 gives all suite-maximum intervals.

A training-fold audit finds five reference tasks that exceed a fitted tool-count maximum. E9 addresses this calibration gap with training-only templates, donor pools and framing caps on real-agent trajectories; E6/E7 retain their original known-population scope.

\paragraph{Online dispatchers}
E8's stand-in planners diagnose sensitivity to unknown future calls and template overflow (Supplement S13). E9 now evaluates this online interface with real model decisions.

\subsection{E9: real agents and training-only defense calibration}\label{sec:real-agent}
GPT-6.1 Sol (low reasoning effort) plans and responds from the retail policy, public tool schemas and actual sandbox observations. A separate context of the same model simulates the user. The rules arm adds one customer-details read before order access and a read before each order write. Each arm has three repeats per task. All 414 training runs on 69 tasks complete before we freeze a shared count template, donor pools and framing caps. Test customers are disjoint from training customers.

The publication snapshot contains 663 completed runs from the planned 690. The primary comparison uses all 41 test tasks with both arms and all repeats complete: 246 runs across 22 customers; three other completed test runs are retained separately. Inclusion depends on completion at a fixed cutoff, not reward or leakage. The online defense inserts training-derived cover on shadow state using only past calls, then flushes unmet counts. Real calls pass through unchanged. We evaluate four program conditions per prompt, with training-fitted attackers, over 2,640 TLS replay sessions. Supplement S14 specifies the protocol, completion accounting and ablations.

\begin{table*}[t]\centering\footnotesize
\caption{E9: real-agent evaluation with training-only defense calibration. Primary snapshot: 41 held-out tasks, 22 customers, two prompts, three repeats (246 episodes). Network values are maximum balanced accuracy (\%) over the fixed attack suite with customer-cluster 95\% intervals. Program = online shadow cover plus constant framing.}\label{tab:r4-primary}
\begin{tabularx}{\textwidth}{@{}Yrrrrr@{}}\toprule
Condition & Cancellation BA [CI] & Return BA [CI] & Task reward (\%) & Bytes$\times$ & Calls$\times$\\\midrule
Baseline & 96.7 [88.9, 100.0] & 93.4 [83.7, 100.0] & 82.1 & 1.00 & 1.00\\
Rules prompt & 96.7 [88.9, 100.0] & 92.8 [83.0, 99.0] & 78.0 & 1.00 & 1.00\\
Program defense & 50.0 [50.0, 64.5] & 50.0 [50.0, 63.1] & 82.1 & 15.33 & 3.83\\
Rules + program & 50.0 [50.0, 66.5] & 64.1 [54.8, 76.1] & 78.0 & 15.70 & 3.90\\
\bottomrule\end{tabularx}
\par\smallskip\raggedright Ratios use raw transport within the same prompt. Program conditions replay the same trajectory, so reward is unchanged by construction. The measurement is TLS~1.2 loopback replay of real model-generated tool calls. The 663-run completion snapshot includes 414 training runs and three further test runs outside the complete-task primary analysis; the planned cohort is 690. Full methods and ablations: Supplement S14.
\end{table*}


Table~\ref{tab:r4-primary} shows strong intent leakage from actual agent behavior: baseline network BA is 96.7\% for cancellation and 93.4\% for return. Rules yield 96.7\% measured workflow compliance (baseline 87.0\%), yet leakage remains 96.7\%/92.8\%. Task reward is 82.1\%/78.0\%; the paired rules-minus-baseline difference is $-4.1$ points (95\% interval $[-10.9,1.7]$). Workflow compliance alone therefore does not deliver metadata privacy in this workload.

With baseline trajectories, training-calibrated online cover plus constant framing reduces both network point estimates to 50.0\%, at 15.3$\times$ bytes and 3.8$\times$ calls. This is a substantial measured reduction under a deployment-relevant calibration constraint. Program replay preserves each trajectory and its reward by construction; it does not measure feedback or latency effects on the agent. Adding rules yields return BA 64.1\%; its paired increase over program-only is 14.1 points with interval $[-4.4,25.7]$, so the interaction remains uncertain. No E9 condition meets the finite-suite criterion for both targets. The result establishes the benefit and measured cost of program-level shaping while showing why prompt behavior, transport and utility require separate evaluation.

\subsection{What the cases establish}
The experiments connect concrete failure modes to the audit's comparison obligations. E5 isolates emitter-dependent cost and count leakage. E6/E7 show that adaptive fitting and encrypted size decoding expose residual tool-intent signals. E9 extends the evidence to real agents and training-only defense calibration: prompt compliance can coexist with high leakage, while online shaping substantially reduces the measured signal at a quantified cost. The cases support these mechanism-level conclusions; they do not estimate the audit's accuracy against independent positive and negative cases.

\soktakeaway{SoK Takeaway 3 (Bind the projection the emitter produces)}{Padding applies to the record sequence it equalizes. Buffered caps, per-event padding and batching have distinct leakage and cost objects. Evaluate the attacker on defended data and price the same emitter that supplies the protection claim.}

\soktakeaway{SoK Takeaway 4 (Evaluate workflow, state and transport together)}{Real-agent workflow rules leave intent inference above 92\%. Training-calibrated online cover with constant framing lowers both baseline network point estimates to 50\%, at 15.3$\times$ bytes and 3.8$\times$ calls. State isolation, task reward and residual inference require distinct measurements.}
